%==========================================================
%  CHAPTER 3: Canada's Housing Crisis — Markets, Affordability,
%              and Policy
%  ECON 204 — Contemporary Economic Issues: A Canadian Perspective
%==========================================================

\chapter{Canada's Housing Crisis: Markets, Affordability, and Policy}
\label{ch:housing}

%----------------------------------------------------------
%  OPENING VIGNETTE
%----------------------------------------------------------

\noindent In February 2022, a detached house in East Vancouver sold for
\$2.1 million. The winning offer was \$400{,}000 above the asking price,
waived all conditions, and was accepted over 22 competing bids. The
buyers --- two engineers in their mid-thirties with combined household
income of approximately \$220{,}000 --- put down \$420{,}000 (20\%) and
took on a mortgage of \$1{,}680{,}000. At that month's prevailing
5-year fixed rate of 3.2\%, their monthly payment was \$8{,}067.

They were the lucky ones. They could afford it.

Across the city, a hospital nurse earning \$72{,}000 per year and a
teacher earning \$65{,}000 per year were renting a two-bedroom apartment
for \$2{,}800 per month --- 40\% of their combined take-home pay. They
had been saving for a down payment for six years, but house prices had
risen faster than their savings every single year. The down payment
they needed was now larger than the house price they had originally
targeted.

And in Prince George --- 800 kilometres to the north --- a single father
working at the mill was paying \$1{,}350 per month for a two-bedroom
apartment in a town where the average household income was \$78{,}000.
He was not in a ``hot market.'' But the housing shortage had spread.

These three households inhabit the same economy, regulated by the same
Bank of Canada, subject to the same federal tax law, and served by the
same CMHC mortgage insurance program. Yet their housing experiences are
radically different --- determined by the intersection of global capital
flows, local zoning bylaws, demographic waves, and macroeconomic policy.

Understanding why Canadian housing became so unaffordable --- and what
could actually fix it --- is one of the most important applied questions
in contemporary Canadian economics.

%----------------------------------------------------------
%  LEARNING OBJECTIVES
%----------------------------------------------------------

\begin{learningobjectives}
  \item Apply the supply-and-demand framework to housing markets,
        explain why the short-run housing supply curve is nearly vertical,
        and predict the price effects of demand shifts.
  \item Derive the user cost of housing formula and use it to evaluate
        the rent-versus-buy decision under different interest rate and
        price appreciation assumptions.
  \item Calculate standard housing affordability measures (price-to-income
        ratio, shelter cost burden, CMHC affordability index) and apply
        the standard mortgage payment formula to scenarios with different
        interest rates.
  \item Explain the supply-side constraints on Canadian housing
        production --- including zoning, land costs, construction labour,
        and permitting --- and evaluate their relative importance.
  \item Analyse the demand-side forces driving Canadian house prices,
        including immigration, financialization, low interest rates, and
        demographic trends.
  \item Describe Canada's rental market dynamics, define the rental
        vacancy rate, and evaluate the economic arguments for and against
        rent control.
  \item Evaluate major federal and provincial housing policy instruments
        using supply-and-demand analysis, and identify which interventions
        address the root cause of unaffordability.
\end{learningobjectives}

%==========================================================
\section{Housing as an Economic Good}
\label{sec:housing_economic_good}
%==========================================================

Housing is an unusual commodity. It is simultaneously a \textbf{consumption
good} (shelter is a basic human need), a \textbf{durable asset} (houses last
50--100 years), a \textbf{financial investment} (households hold most of their
net worth in housing), and a \textbf{location-specific good} (a house in
Vancouver cannot substitute for a house in Saskatoon).

These characteristics give housing markets properties that distinguish them
from markets for ordinary goods:

\begin{itemize}
  \item \textbf{Highly inelastic short-run supply.} It takes 18--36 months
        to plan, permit, and build new housing. The existing stock of homes
        changes slowly. A sudden increase in demand therefore causes prices
        to rise sharply in the short run, even if the long-run supply
        response would eventually moderate prices.

  \item \textbf{Spatially fixed.} Housing supply cannot respond to demand
        by shipping inventory from low-demand to high-demand regions. Land
        in downtown Toronto is permanently scarce in a way that land in
        Kansas is not. This creates permanent scarcity rents in desirable
        locations that compound over time.

  \item \textbf{Heterogeneous.} No two homes are identical. Differences
        in location, size, vintage, condition, and neighbourhood make
        ``the housing market'' really an overlapping set of segmented
        submarkets that interact but behave differently.

  \item \textbf{Highly leveraged.} Most housing purchases are debt-financed.
        A household making a 10\% down payment is borrowing 90 cents of
        every dollar of housing value. This leverage amplifies both gains
        and losses, makes housing demand extremely sensitive to interest
        rates, and connects housing markets tightly to the credit system.

  \item \textbf{Regulated at every level.} Zoning bylaws (municipal),
        building codes (provincial), mortgage insurance rules (federal),
        rent regulations (provincial), environmental assessments
        (federal/provincial), and development charges (municipal) all
        constrain the supply of, and demand for, housing simultaneously.
\end{itemize}

\begin{defbox}{The Rental Market and the Ownership Market}
The housing market has two interconnected segments.

The \textbf{ownership market} prices houses as assets. The equilibrium
price reflects the expected present value of future housing services,
including the option value of residing in a particular location.

The \textbf{rental market} prices housing as a flow of services per unit
of time. Renters pay for the right to occupy a dwelling without acquiring
the asset.

The two markets are linked: when owning becomes more expensive relative
to renting (e.g.\ because mortgage rates rise), demand shifts from
ownership to rental, pushing rental prices up. Conversely, when rents
are very high, the yield on owning rental property rises, incentivising
investors to buy.

In Canada's current crisis, both markets have tightened simultaneously ---
a phenomenon explained by the coincidence of record immigration, low rental
vacancy, high mortgage rates, and constrained new supply.
\end{defbox}

%==========================================================
\section{Supply and Demand in Housing Markets}
\label{sec:housing_supply_demand}
%==========================================================

%----------------------------------------------------------
\subsection{The Demand for Housing}
%----------------------------------------------------------

The demand for housing is driven by:

\begin{itemize}
  \item \textbf{Population growth and household formation:} More people
        require more dwellings. Canada's record immigration flows
        (431{,}645 permanent residents in 2022; over 1 million people
        added to the population including non-permanent residents) have
        been the single largest driver of housing demand in the mid-2020s.

  \item \textbf{Income:} Higher household income increases both the
        ability and willingness to pay for housing. Housing is a
        \term{normal good} with an income elasticity of demand
        estimated between 0.5 and 1.0.

  \item \textbf{Interest rates:} The monthly cost of servicing a mortgage
        depends directly on the interest rate. Lower rates reduce the
        flow cost of owning a given quantity of housing, increasing demand.
        The 2020--2021 rate cuts to 0.25\% dramatically reduced mortgage
        payments and turbocharged demand.

  \item \textbf{Expected capital gains:} If buyers expect house prices
        to rise, the expected capital gain reduces the user cost of owning
        and shifts demand rightward. Extrapolative expectations can generate
        self-fulfilling price bubbles.

  \item \textbf{Investor demand:} A significant share of housing demand
        comes from investors who do not occupy the units. In some markets,
        investor demand competes directly with owner-occupiers and pushes
        prices above what the fundamental rental income stream would justify.
\end{itemize}

%----------------------------------------------------------
\subsection{The Supply of Housing}
%----------------------------------------------------------

The supply of housing in any given market is determined by:

\begin{itemize}
  \item \textbf{Existing stock:} The vast majority of housing supply in
        any year is the existing stock of homes. In Canada, approximately
        15 million private dwellings are occupied. New construction in a
        typical year (200{,}000--260{,}000 units) adds only 1.3--1.7\%
        to the stock.

  \item \textbf{New construction:} \term{Housing starts} measure the number
        of new residential units on which construction has begun. Their
        determinants include land prices, labour costs, material costs,
        developer financing costs, development charges, permitting timelines,
        and zoning regulations.

  \item \textbf{Conversion and densification:} Basement suites, laneway
        houses, and conversions of single-family homes to multiplexes
        can increase supply without requiring greenfield land.

  \item \textbf{Demolition and obsolescence:} A small fraction of the
        existing stock is demolished or becomes uninhabitable each year.
        Net new supply = New construction $-$ Demolitions.
\end{itemize}

%----------------------------------------------------------
\subsection{Supply Elasticity and Its Consequences}
%----------------------------------------------------------

The \term{housing supply elasticity} measures how responsive new construction
is to price increases:

\begin{equation}
  \varepsilon_s = \frac{\%\Delta Q_s^{\text{housing}}}{\%\Delta P_{\text{housing}}}
  \label{eq:supply_elasticity}
\end{equation}

Studies for major Canadian cities consistently find supply elasticities
well below 1.0 --- typically between 0.2 and 0.6 for the short run.
This means that a 10\% increase in house prices produces only a 2--6\%
increase in housing supply. The reasons for this inelasticity are explored
in Section~\ref{sec:supply_constraints}.

Figure~\ref{fig:housing_sd} illustrates the consequences of low supply
elasticity for housing price dynamics.

\begin{figure}[H]
\centering
\begin{tikzpicture}[
  axisline/.style={-Stealth, thick, color=gray!70},
  lbl/.style={font=\small}
]

%---- PANEL A: INELASTIC SUPPLY ----
\begin{scope}[xshift=0cm]
\node[font=\small\bfseries, align=center] at (3.2,5.8)
  {Panel A: Inelastic Supply (Canada)};
\draw[axisline] (0,0) -- (0,5.5) node[above, lbl] {$P$};
\draw[axisline] (0,0) -- (6.0,0) node[right, lbl] {$Q$};
% Inelastic supply
\draw[policyblue, thick] (2.5,0.3) -- (3.0,5.0)
  node[right, font=\footnotesize, color=policyblue] {$S_{\text{inelastic}}$};
% D1
\draw[unbcgreen, thick] (0.5,4.5) -- (5.5,0.3)
  node[right, font=\footnotesize, color=unbcgreen] {$D_1$};
% D2
\draw[casered, thick, dashed] (1.5,4.8) -- (6.0,0.7)
  node[right, font=\footnotesize, color=casered] {$D_2$};
\fill[unbcgreen] (2.68,2.2) circle (2.8pt);
\fill[casered]   (2.83,3.9) circle (2.8pt);
\draw[dashed, gray!50, thin] (0,2.2) -- (2.68,2.2);
\draw[dashed, gray!50, thin] (0,3.9) -- (2.83,3.9);
\draw[dashed, gray!50, thin] (2.68,0) -- (2.68,2.2);
\draw[dashed, gray!50, thin] (2.83,0) -- (2.83,3.9);
\node[left, font=\footnotesize] at (0,2.2) {$P_1$};
\node[left, font=\footnotesize] at (0,3.9) {$P_2^A$};
\node[below, font=\footnotesize] at (2.68,0) {$Q_1$};
\node[below, font=\footnotesize] at (2.83,0) {$Q_2^A$};
\draw[<->, casered!70, thick] (0.5,2.2) -- (0.5,3.9)
  node[midway, left, font=\footnotesize, color=casered] {$\Delta P$ large};
\draw[<->, policyblue!70, thick] (2.68,-0.5) -- (2.83,-0.5)
  node[midway, below, font=\footnotesize, color=policyblue] {$\Delta Q$ small};
\end{scope}

%---- PANEL B: ELASTIC SUPPLY ----
\begin{scope}[xshift=7.5cm]
\node[font=\small\bfseries, align=center] at (3.2,5.8)
  {Panel B: Elastic Supply (Policy Goal)};
\draw[axisline] (0,0) -- (0,5.5) node[above, lbl] {$P$};
\draw[axisline] (0,0) -- (6.0,0) node[right, lbl] {$Q$};
\draw[policyblue, thick] (0.5,1.8) -- (5.5,2.8)
  node[right, font=\footnotesize, color=policyblue] {$S_{\text{elastic}}$};
\draw[unbcgreen, thick] (0.5,4.5) -- (5.5,0.3)
  node[right, font=\footnotesize, color=unbcgreen] {$D_1$};
\draw[casered, thick, dashed] (1.5,4.8) -- (6.0,0.7)
  node[right, font=\footnotesize, color=casered] {$D_2$};
\fill[unbcgreen] (2.9,2.12) circle (2.8pt);
\fill[casered]   (4.15,2.38) circle (2.8pt);
\draw[dashed, gray!50, thin] (0,2.12) -- (2.9,2.12);
\draw[dashed, gray!50, thin] (0,2.38) -- (4.15,2.38);
\draw[dashed, gray!50, thin] (2.9,0) -- (2.9,2.12);
\draw[dashed, gray!50, thin] (4.15,0) -- (4.15,2.38);
\node[left, font=\footnotesize] at (0,2.12) {$P_1$};
\node[left, font=\footnotesize] at (0,2.38) {$P_2^B$};
\node[below, font=\footnotesize] at (2.9,0) {$Q_1$};
\node[below, font=\footnotesize] at (4.15,0) {$Q_2^B$};
\draw[<->, casered!70, thick] (0.5,2.12) -- (0.5,2.38)
  node[midway, left, font=\footnotesize, color=casered] {$\Delta P$ small};
\draw[<->, policyblue!70, thick] (2.9,-0.5) -- (4.15,-0.5)
  node[midway, below, font=\footnotesize, color=policyblue] {$\Delta Q$ large};
\end{scope}
\end{tikzpicture}
\caption{Housing Supply Elasticity and Price Dynamics.
\textbf{Panel A} shows Canada's current reality: the supply curve is nearly
vertical (inelastic) because zoning, land scarcity, and slow permitting
prevent new housing from responding quickly to demand. When demand shifts
from $D_1$ to $D_2$, prices rise sharply while quantity barely changes.
\textbf{Panel B} shows the policy goal: a more elastic supply curve absorbs
the same demand shift with a much smaller price increase and a much larger
quantity response. The fundamental insight is that housing affordability is
primarily a supply problem, and the supply curve's slope is the critical
policy parameter.}
\label{fig:housing_sd}
\end{figure}

%==========================================================
\section{The User Cost of Housing: Buy or Rent?}
\label{sec:user_cost}
%==========================================================

The decision to buy or rent a home is analysed using the \term{user cost
of housing} framework. The user cost is the annual economic cost of occupying
a home you own, expressed as a fraction of its market value:

\begin{equation}
  \text{UC} = (r + \delta + \tau - g) \times P
  \label{eq:user_cost}
\end{equation}

\noindent where:
\begin{itemize}
  \item $r$ = real mortgage interest rate
  \item $\delta$ = annual depreciation rate (approximately 1.5--2.5\%)
  \item $\tau$ = property tax rate (approximately 0.5--1.5\%)
  \item $g$ = expected real capital gains rate (annual price appreciation above inflation)
  \item $P$ = current price of the home
\end{itemize}

Renting is more economical when the annual rent is less than UC;
buying is more economical when annual rent exceeds UC.

\noindent\textbf{Worked Example 3.1 --- User Cost Under Different Conditions}

\noindent Consider a home priced at \$700{,}000 in Ottawa under three scenarios:

\begin{center}
\renewcommand{\arraystretch}{1.35}
\begin{tabular}{lrrr}
\toprule
\textbf{Parameter} & \textbf{2021 (Low Rate)} & \textbf{2023 (High Rate)}
                   & \textbf{2026 (Normal)} \\
\midrule
Mortgage rate $r$ (\%)         & 2.0 & 5.8 & 4.5 \\
Depreciation $\delta$ (\%)     & 2.0 & 2.0 & 2.0 \\
Property tax $\tau$ (\%)       & 1.0 & 1.0 & 1.0 \\
Expected real gain $g$ (\%)    & 4.0 & 0.0 & 1.5 \\
\midrule
\textbf{User cost rate (\%)}   & \textbf{1.0} & \textbf{8.8} & \textbf{6.0} \\
\textbf{Annual user cost (\$)} & \textbf{\$7{,}000} & \textbf{\$61{,}600}
                                & \textbf{\$42{,}000} \\
\textbf{Monthly equivalent}    & \textbf{\$583} & \textbf{\$5{,}133}
                                & \textbf{\$3{,}500} \\
\bottomrule
\end{tabular}
\end{center}

\noindent\textit{Interpretation:} In 2021, the monthly user cost of owning
a \$700{,}000 home was only \$583 --- far below any comparable rent.
This explains the demand surge of 2020--2022: owning was far cheaper than
renting in user cost terms. By 2023, the monthly user cost had risen to
\$5{,}133, making owning economically irrational relative to renting
in many markets.

\begin{thinkbox}
Why did house prices continue rising in 2021 even though prices were already
``unaffordably high'' by price-to-income measures? The user cost framework
answers this. When $g > r$ (expected capital gains exceed the interest rate),
the user cost can approach zero or become negative. The more expensive the
house, the more capital gain the owner earns --- so higher prices make buying
\textit{more} attractive, not less. This creates a feedback loop that sustains
price bubbles: rising prices increase expected future gains, reduce user cost,
attract more buyers, push prices higher, and so on. The loop breaks when
interest rates rise sharply (as in 2022), expected capital gains collapse,
and user costs suddenly jump.
\end{thinkbox}

%==========================================================
\section{Measuring Housing Affordability}
\label{sec:affordability_measurement}
%==========================================================

%----------------------------------------------------------
\subsection{The Price-to-Income Ratio}
%----------------------------------------------------------

The \term{price-to-income ratio (PTI)} is the most widely reported
affordability benchmark:

\begin{equation}
  \text{PTI} = \frac{\text{Median House Price}}{\text{Median Annual Household Income}}
  \label{eq:pti}
\end{equation}

A PTI of 3.0 is traditionally considered ``affordable.''
Above 5.0 is ``seriously unaffordable.'' Above 9.0 is ``severely unaffordable.''

\begin{figure}[H]
\centering
\begin{tikzpicture}
\begin{axis}[
  ybar,
  bar width=0.48cm,
  width=13cm, height=7cm,
  symbolic x coords={{Halifax},{Winnipeg},{Calgary},{Montreal},{Ottawa},
                     {CAN Avg},{Houston},{Berlin},{Toronto},{NYC},{Sydney},{Vancouver}},
  xtick=data,
  x tick label style={rotate=40, anchor=east, font=\small},
  ylabel={Price-to-Income Ratio, 2024},
  ylabel style={font=\small},
  ymin=0, ymax=15,
  ytick={0,2,4,6,8,10,12,14},
  yticklabel style={font=\small},
  ymajorgrids=true,
  grid style={dashed, gray!25},
  enlarge x limits=0.05,
  nodes near coords,
  nodes near coords style={font=\footnotesize},
  every axis plot/.append style={fill=policyblue!60, draw=policyblue!80}
]
\addplot coordinates {
  ({Halifax},    6.8)
  ({Winnipeg},   4.8)
  ({Calgary},    6.5)
  ({Montreal},   6.2)
  ({Ottawa},     7.2)
  ({CAN Avg},    8.1)
  ({Houston},    3.8)
  ({Berlin},     8.9)
  ({Toronto},   10.2)
  ({NYC},        9.8)
  ({Sydney},    13.3)
  ({Vancouver}, 12.8)
};
\addplot[unbcgreen, dashed, thick, forget plot]
  coordinates {({Halifax},3.0)({Vancouver},3.0)};
\addplot[casered, dashed, thick, forget plot]
  coordinates {({Halifax},5.0)({Vancouver},5.0)};
\node[font=\footnotesize, unbcgreen] at (axis cs:{Winnipeg},3.45)
  {``Affordable'' (3.0)};
\node[font=\footnotesize, casered] at (axis cs:{Winnipeg},5.45)
  {``Unaffordable'' (5.0)};
\end{axis}
\end{tikzpicture}
\caption{Price-to-Income Ratio, Selected Canadian and International Cities, 2024.
Vancouver (12.8) and Toronto (10.2) rank among the least affordable cities
globally. Only Winnipeg (4.8) of the major Canadian cities approaches the
historic ``affordable'' norm of 3.0 (green dashed line). The international
comparison is instructive: Houston (3.8) is a major North American city with
permissive land-use rules; Berlin (8.9) has high prices but also a large
social housing sector that insulates lower-income households from full market
exposure. Canada's affordability problem is real, severe, and primarily
concentrated in its two largest metropolitan areas.}
\label{fig:pti_comparison}
\source{Demographia International Housing Affordability Survey, 2024;
CMHC Housing Market Outlook, 2024.}
\end{figure}

%----------------------------------------------------------
\subsection{The Mortgage Payment Formula}
%----------------------------------------------------------

A second affordability measure is the \term{shelter cost burden}: the fraction
of income required to service housing costs. The conventional threshold is
30\% of gross income.

The standard Canadian mortgage payment formula is:

\begin{equation}
  M = P_0 \times \frac{r_m (1 + r_m)^n}{(1 + r_m)^n - 1}
  \label{eq:mortgage_payment}
\end{equation}

\noindent where $M$ is the monthly payment, $P_0$ is the principal (price
minus down payment), $r_m$ is the effective monthly rate derived from the
annual rate under the Canadian semi-annual compounding convention
($r_m = (1 + r_{\text{annual}}/2)^{1/6} - 1$), and $n$ is total monthly
payments (amortization years $\times 12$).

\begin{table}[H]
\centering
\caption{Monthly Mortgage Payment and Affordability, \$700{,}000 Home,
20\% Down (\$140{,}000), \$560{,}000 Principal, 25-Year Amortization}
\label{tab:mortgage_payments}
\renewcommand{\arraystretch}{1.35}
\begin{tabular}{lrrrr}
\toprule
\textbf{Rate (\%)} & \textbf{Monthly Payment} & \textbf{Annual Cost}
                   & \textbf{Income at 30\% Rule} & \textbf{Total Interest} \\
\midrule
2.0 & \$2{,}370 & \$28{,}440 & \$94{,}800  & \$150{,}800 \\
3.0 & \$2{,}646 & \$31{,}752 & \$105{,}840 & \$233{,}600 \\
4.0 & \$2{,}936 & \$35{,}232 & \$117{,}440 & \$320{,}800 \\
5.0 & \$3{,}263 & \$39{,}156 & \$130{,}520 & \$418{,}900 \\
5.8 & \$3{,}531 & \$42{,}372 & \$141{,}240 & \$499{,}300 \\
7.0 & \$3{,}921 & \$47{,}052 & \$156{,}840 & \$616{,}600 \\
\bottomrule
\end{tabular}
\source{Author calculations using Equation~\ref{eq:mortgage_payment}
under Canadian semi-annual compounding. ``Income at 30\% rule'' is the
annual gross income required for the mortgage payment to represent exactly
30\% of income. Total interest assumes full 25-year amortization.}
\end{table}

At 5.8\% --- the approximate 5-year fixed rate prevailing through 2023 ---
a household buying a \$700{,}000 home needed gross income of at least
\$141{,}240 per year to meet the 30\% threshold. In Vancouver, where the
benchmark price exceeded \$1.1 million through 2024, the required income
at 5.8\% exceeded \$225{,}000.

%----------------------------------------------------------
\subsection{The CMHC Affordability Index}
%----------------------------------------------------------

CMHC publishes a \term{Housing Affordability Index} measuring mortgage
payment, property taxes, and utilities as a percentage of median household
income for a median-priced home.

\begin{table}[H]
\centering
\caption{CMHC Housing Affordability Index, Selected Cities, 2015--2024}
\label{tab:affordability_index}
\renewcommand{\arraystretch}{1.30}
\begin{tabular}{lrrrrrr}
\toprule
\textbf{City} & \textbf{2015} & \textbf{2018} & \textbf{2020}
              & \textbf{2022} & \textbf{2023} & \textbf{2024} \\
\midrule
Vancouver  & 80.2 & 87.6 & 61.4 & 93.8 & 90.2 & 92.1 \\
Toronto    & 56.8 & 74.2 & 67.3 & 89.4 & 82.7 & 75.3 \\
Ottawa     & 38.2 & 42.1 & 44.5 & 62.3 & 58.7 & 48.2 \\
Calgary    & 35.6 & 32.1 & 31.8 & 42.5 & 45.1 & 38.4 \\
Montr\'{e}al & 34.0 & 38.4 & 41.2 & 55.2 & 49.8 & 42.3 \\
Halifax    & 26.1 & 28.3 & 30.2 & 49.3 & 46.5 & 44.2 \\
Winnipeg   & 25.8 & 26.5 & 27.3 & 33.2 & 31.5 & 29.4 \\
\textbf{National avg} & \textbf{40.2} & \textbf{50.1} & \textbf{46.8}
                      & \textbf{65.3} & \textbf{58.4} & \textbf{53.2} \\
\bottomrule
\end{tabular}
\source{Canada Mortgage and Housing Corporation, Housing Market Outlook,
2024. Index = (mortgage payment + taxes + utilities) as \% of median
household income. Values above 30\% indicate unaffordable housing.}
\end{table}

Vancouver's 2024 index of 92.1\% means the median Vancouver household
buying the median-priced home would spend 92\% of gross income on housing
costs --- mathematically impossible for most families. The national average
peaked at 65.3\% in 2022 and has partially improved since, but remains
far above the pre-crisis levels of the early 2010s.

%==========================================================
\section{The Historical Price Record}
\label{sec:house_price_history}
%==========================================================

\begin{figure}[H]
\centering
\begin{tikzpicture}
\begin{axis}[
  width=14cm, height=7.5cm,
  xlabel={Year},
  ylabel={MLS Benchmark Price (\$CAD thousands)},
  xlabel style={font=\small},
  ylabel style={font=\small},
  xmin=1999, xmax=2027,
  ymin=0, ymax=1600,
  xtick={2000,2004,2008,2012,2016,2020,2024},
  xticklabel style={font=\small},
  ytick={0,200,400,600,800,1000,1200,1400,1600},
  yticklabel style={font=\small},
  ymajorgrids=true,
  grid style={dashed, gray!20},
  legend pos=north west,
  legend style={font=\small},
  clip=false
]
\addplot[gray!70, thick] coordinates {
  (2000,163)(2004,233)(2008,302)(2010,340)(2012,365)(2014,400)
  (2016,480)(2018,488)(2020,530)(2021,716)(2022,796)
  (2023,650)(2024,700)(2025,710)(2026,720)
};
\addlegendentry{Canada national}
\addplot[policyblue, thick] coordinates {
  (2000,270)(2004,370)(2008,520)(2010,598)(2012,683)(2014,720)
  (2016,1050)(2018,870)(2020,1004)(2021,1200)(2022,1370)
  (2023,1170)(2024,1250)(2025,1270)(2026,1290)
};
\addlegendentry{Vancouver}
\addplot[casered, thick, dashed] coordinates {
  (2000,250)(2004,310)(2008,362)(2010,431)(2012,497)(2014,566)
  (2016,730)(2018,736)(2020,930)(2021,1163)(2022,1340)
  (2023,1100)(2024,1148)(2025,1160)(2026,1175)
};
\addlegendentry{Toronto}
\addplot[dataorange, thick, dotted] coordinates {
  (2000,160)(2006,340)(2008,380)(2010,365)(2016,420)
  (2020,420)(2021,475)(2022,560)(2023,570)(2024,600)(2025,615)(2026,625)
};
\addlegendentry{Calgary}
\node[font=\footnotesize, align=center, color=gray] at (axis cs:2008.5,160)
  {GFC};
\node[font=\footnotesize, color=unbcgreen, align=left] at (axis cs:2020.2,420)
  {COVID\\low rates};
\node[font=\footnotesize, color=casered, align=right] at (axis cs:2022.5,1450)
  {2022 peak};
\end{axis}
\end{tikzpicture}
\caption{MLS Composite Benchmark House Prices, Canada and Selected Cities,
2000--2026. Canadian house prices have risen more than four-fold in nominal
terms since 2000; Vancouver and Toronto five- to six-fold. Key episodes:
the 2008 Global Financial Crisis caused only a brief pause; the 2016 BC
Foreign Buyer Tax caused a temporary Vancouver correction; COVID-era
ultra-low rates (2020--2021) triggered the most explosive price surge in
history; and the 2022--2023 rate hikes caused the first significant correction,
though prices remain far above pre-pandemic levels. Calgary's flatter profile
reflects oil price cycle sensitivity.}
\label{fig:house_prices}
\source{Canadian Real Estate Association (CREA), MLS Home Price Index.
2025--2026 estimated.}
\end{figure}

%==========================================================
\section{Supply-Side Constraints on Canadian Housing}
\label{sec:supply_constraints}
%==========================================================

The central analytical question in housing economics is whether price
increases are primarily driven by demand or supply. The academic evidence
increasingly points to \textbf{supply constraints as the primary long-run
driver}. The same demand pressures exist in many countries without producing
the same price outcomes --- what differs is how quickly supply responds.

%----------------------------------------------------------
\subsection{Zoning and Land Use Regulation}
%----------------------------------------------------------

The most important supply-side constraint is \term{zoning} --- the municipal
regulation of how land may be used. In most Canadian cities, the majority
of residential land is zoned exclusively for single-family detached housing.
This ``Euclidean'' zoning model prevents the construction of apartments,
townhouses, or duplexes on most residential land regardless of market demand.

The economic consequences are profound:

\begin{itemize}
  \item \textbf{Land value inflation:} When zoning prevents densification,
        the value of a single-family lot captures the entire surplus from
        excluding all alternative, potentially more productive uses.

  \item \textbf{Reduced supply elasticity:} Zoning directly reduces
        $\varepsilon_s$ (Equation~\ref{eq:supply_elasticity}). Developers
        who want to build apartments cannot do so without lengthy, costly,
        and uncertain rezoning processes.

  \item \textbf{Parking minimums and setback requirements:} Many municipalities
        require 1--2 parking spaces per unit, minimum lot sizes, setbacks,
        and maximum lot coverage. Each constraint reduces the number of
        units buildable per parcel.
\end{itemize}

\begin{canexample}{British Columbia's 2023 Zoning Reform}
In December 2023, BC passed legislation requiring all municipalities in
Metro Vancouver and most municipalities with populations over 5{,}000 to
permit ``small scale multi-unit housing'' --- duplexes, triplexes,
fourplexes, and sixplexes --- on lots previously zoned exclusively for
single-family homes. This \term{upzoning} legislation was the most
significant provincial override of municipal zoning in Canadian history.

Economists estimate the reform could permit hundreds of thousands of
additional units over the next decade if complemented by appropriate
infrastructure investment and streamlined permitting. Early evidence
from 2025 suggests modest but positive effects on building permit
applications in affected areas.
\end{canexample}

%----------------------------------------------------------
\subsection{Development Costs and Construction Capacity}
%----------------------------------------------------------

Even where zoning permits higher density, additional cost factors constrain supply:

\begin{itemize}
  \item \textbf{Development charges:} In Toronto, development charges
        on a new condo unit reached approximately \$150{,}000--\$200{,}000
        in 2023 --- representing 15--25\% of the unit's total cost, ultimately
        borne by buyers through higher prices.

  \item \textbf{Construction labour shortages:} Canada's construction sector
        faces significant skilled labour shortages --- carpenters, electricians,
        plumbers, ironworkers. Construction wages have risen faster than
        overall wages since 2020.

  \item \textbf{Permitting timelines:} The time from a development application
        to a building permit in major Canadian cities averages 12--36 months.
        Each month of delay costs developers carrying charges on their land.

  \item \textbf{High-rise construction costs:} Tall concrete and steel
        construction costs \$350--\$600 per square foot in Canada, compared
        to \$150--\$250 for wood-frame low-rise. The economics only work
        at very high price levels.
\end{itemize}

%----------------------------------------------------------
\subsection{The Supply-Demand Gap: Housing Starts vs. Population Growth}
%----------------------------------------------------------

\begin{figure}[H]
\centering
\begin{tikzpicture}
\begin{axis}[
  name=ax1,
  width=14cm, height=7cm,
  xlabel={Year},
  ylabel={Housing Starts (thousands of units)},
  ylabel style={font=\small, color=policyblue},
  xlabel style={font=\small},
  xmin=1999, xmax=2025.5,
  ymin=100, ymax=320,
  xtick={2000,2004,2008,2012,2016,2020,2024},
  xticklabel style={font=\small},
  ytick={100,150,200,250,300},
  yticklabel style={font=\small, color=policyblue},
  axis y line*=left,
  axis x line*=bottom,
  ymajorgrids=true,
  grid style={dashed, gray!20},
  clip=false,
  legend style={font=\small, at={(0.02,0.98)}, anchor=north west}
]
\addplot[policyblue, thick, mark=*, mark size=1.5pt] coordinates {
  (2000,151)(2002,205)(2004,233)(2006,228)(2008,211)(2009,149)
  (2010,190)(2012,215)(2014,190)(2016,198)(2018,212)(2020,218)
  (2021,271)(2022,263)(2023,239)(2024,245)
};
\addlegendentry{Housing Starts (left)}
\node[font=\footnotesize, policyblue, align=center] at (axis cs:2022.5,295)
  {Need $\approx$450K+ units/yr};
\draw[->, policyblue!60, thin] (axis cs:2022,283) -- (axis cs:2022,265);
\end{axis}
\begin{axis}[
  width=14cm, height=7cm,
  ylabel={Annual Population Growth (thousands)},
  ylabel style={font=\small, color=casered},
  xmin=1999, xmax=2025.5,
  ymin=200, ymax=1400,
  ytick={200,400,600,800,1000,1200,1400},
  yticklabel style={font=\small, color=casered},
  axis y line*=right,
  axis x line=none,
  legend style={font=\small, at={(0.98,0.98)}, anchor=north east},
  clip=false
]
\addplot[casered, thick, dashed, mark=square*, mark size=1.5pt] coordinates {
  (2000,362)(2002,339)(2004,337)(2006,355)(2008,393)(2010,338)
  (2012,357)(2014,388)(2016,339)(2018,528)(2020,329)(2021,507)
  (2022,1053)(2023,1159)(2024,1200)
};
\addlegendentry{Population Growth (right)}
\end{axis}
\end{tikzpicture}
\caption{Housing Starts vs.\ Annual Population Growth, Canada, 2000--2024.
Until approximately 2017, housing starts and population growth moved in rough
proportion. From 2022 onwards, Canada's record immigration levels produced
population growth of over 1 million persons per year --- requiring approximately
400{,}000--480{,}000 new housing units annually (at 2.5 persons per household)
--- while actual construction remained at 240{,}000--265{,}000 starts.
This structural supply shortfall is the primary driver of the post-2021
rental vacancy crisis and rent inflation.}
\label{fig:starts_vs_population}
\source{Statistics Canada, Tables 34-10-0035-01 (housing starts) and
17-10-0008-01 (population growth estimates).}
\end{figure}

%==========================================================
\section{Demand-Side Forces: Immigration, Finance, and Demographics}
\label{sec:demand_side}
%==========================================================

While supply constraints are the primary structural explanation for high
prices, several demand-side forces have amplified and accelerated the
affordability crisis.

%----------------------------------------------------------
\subsection{Immigration and Population Growth}
%----------------------------------------------------------

Canada's record immigration of 431{,}645 permanent residents in 2022 (and
approximately 1 million total population additions including non-permanent
residents) represents the single largest year-over-year demand shock to
housing in Canadian history. The housing consequences are significant because:

\begin{enumerate}
  \item Newcomers must immediately find accommodation.
  \item Newcomers concentrate geographically: 57\% of 2022 permanent
        residents settled in Ontario or British Columbia --- precisely the
        markets with lowest vacancy rates.
  \item Non-permanent residents (international students, temporary foreign
        workers) are disproportionately renters, adding pressure to already
        strained rental markets.
\end{enumerate}

Immigration does not make immigration bad policy --- Chapter~\ref{ch:immigration}
presents the strong evidence for its economic benefits. But it does require
that housing supply policy be coordinated with immigration policy. Canada
set record immigration targets without simultaneously implementing the
supply-side reforms needed to house the newcomers. This policy coordination
failure contributed significantly to the 2022--2024 rental crisis.

%----------------------------------------------------------
\subsection{The Interest Rate Channel}
%----------------------------------------------------------

As established in Chapter~\ref{ch:inflation}, the Bank of Canada's decision
to hold the overnight rate at 0.25\% through 2020--2021 created deeply
negative real interest rates. Through the user cost mechanism
(Equation~\ref{eq:user_cost}), this dramatically reduced the annual cost
of owning housing:

\begin{itemize}
  \item At $r = 2\%$ and $g = 5\%$ expected capital gains, the user cost
        rate was approximately $-0.5\%$ to $+1.0\%$.
  \item This meant the annual cost of owning a \$1{,}000{,}000 home was
        \$0--\$10{,}000, while comparable rental costs ran \$36{,}000--\$48{,}000
        per year in Vancouver or Toronto.
  \item The rent-vs-buy calculation was decisively in favour of buying,
        generating enormous ownership demand.
\end{itemize}

\begin{figure}[H]
\centering
\begin{tikzpicture}
\begin{axis}[
  width=14cm, height=6.5cm,
  xlabel={Year},
  ylabel={Rate (\%)},
  xlabel style={font=\small},
  ylabel style={font=\small},
  xmin=2014.5, xmax=2026.5,
  ymin=0, ymax=8,
  xtick={2015,2016,2017,2018,2019,2020,2021,2022,2023,2024,2025,2026},
  xticklabel style={rotate=35, anchor=east, font=\small},
  ytick={0,1,2,3,4,5,6,7,8},
  yticklabel style={font=\small},
  ymajorgrids=true,
  grid style={dashed, gray!20},
  legend pos=north west,
  legend style={font=\small},
  clip=false
]
\addplot[policyblue, thick, mark=*, mark size=1.5pt] coordinates {
  (2015,4.6)(2016,4.6)(2017,4.6)(2018,5.3)(2019,3.5)
  (2020,2.2)(2021,2.0)(2022,5.1)(2023,6.1)(2024,5.2)(2025,4.6)(2026,4.3)
};
\addlegendentry{5-year fixed mortgage rate}
\addplot[casered, thick, dashed, mark=square*, mark size=1.5pt] coordinates {
  (2015,2.7)(2016,2.7)(2017,3.2)(2018,3.9)(2019,3.5)
  (2020,1.5)(2021,1.4)(2022,4.8)(2023,6.9)(2024,5.4)(2025,4.0)(2026,3.6)
};
\addlegendentry{Variable rate (approx.)}
\addplot[unbcgreen!70, dotted, thick] coordinates {
  (2015,0.5)(2016,0.5)(2017,1.0)(2018,1.75)(2019,1.75)
  (2020,0.25)(2021,0.25)(2022,4.25)(2023,5.0)(2024,3.25)(2025,3.0)(2026,2.75)
};
\addlegendentry{Bank of Canada overnight rate}
\node[font=\footnotesize, policyblue, align=center] at (axis cs:2020.5,0.8)
  {Historic\\rate lows};
\node[font=\footnotesize, casered] at (axis cs:2023,7.3)
  {Post-hike peak};
\end{axis}
\end{tikzpicture}
\caption{Mortgage Rates vs.\ Bank of Canada Overnight Rate, 2015--2026.
The 5-year fixed mortgage rate fell to a historic low of approximately 2.0\%
in 2021, driving a massive surge in housing demand. The rate-hike cycle
drove the 5-year fixed rate above 6\% by mid-2023 --- more than tripling
the cost of new mortgage borrowing. Variable-rate holders saw their payments
rise almost instantly, tracking the overnight rate with minimal lag.}
\label{fig:mortgage_rates}
\source{Bank of Canada; Statistics Canada. 2025--2026 estimated.}
\end{figure}

%----------------------------------------------------------
\subsection{Financialization of Housing}
%----------------------------------------------------------

A third demand-side factor is the \term{financialization of housing} ---
the increasing treatment of residential real estate as a financial asset class:

\begin{itemize}
  \item \textbf{Investor purchasing:} CMHC estimates that investors
        accounted for 20--30\% of condominium purchases in Toronto and
        Vancouver in recent years, competing with first-time buyers.

  \item \textbf{Short-term rentals:} Airbnb and similar platforms converted
        approximately 31{,}000 long-term rental units into tourist
        accommodation in Canada's 10 largest cities.

  \item \textbf{Vacancy and speculation:} Some buyers hold units vacant
        as stores of value. Several Canadian cities have introduced vacancy
        taxes (Vancouver Empty Homes Tax, BC Speculation and Vacancy Tax)
        that address the problem directly and generate data to measure its scale.
\end{itemize}

%==========================================================
\section{The Rental Crisis}
\label{sec:rental_crisis}
%==========================================================

As buying became impossible for an increasing share of Canadians, demand
for rental housing surged --- but rental supply grew slowly, pushing vacancy
rates to historic lows and rents to historic highs.

\begin{figure}[H]
\centering
\begin{tikzpicture}
\begin{axis}[
  width=14cm, height=6.5cm,
  xlabel={Year},
  ylabel={Rental Vacancy Rate (\%)},
  xlabel style={font=\small},
  ylabel style={font=\small},
  xmin=2013.5, xmax=2025,
  ymin=0, ymax=7,
  xtick={2014,2015,2016,2017,2018,2019,2020,2021,2022,2023,2024},
  xticklabel style={font=\small},
  ytick={0,1,2,3,4,5,6,7},
  yticklabel style={font=\small},
  ymajorgrids=true,
  grid style={dashed, gray!20},
  legend pos=north east,
  legend style={font=\small},
  clip=false
]
\addplot[gray!50, dotted, thick, forget plot]
  coordinates {(2013.5,3.0)(2025,3.0)};
\node[right, font=\footnotesize, color=gray!70]
  at (axis cs:2024.7,3.0) {};
\node[font=\footnotesize, align=left, color=gray!70]
  at (axis cs:2015.5,3.3) {Healthy (3\%)};
\addplot[policyblue, thick, mark=*, mark size=1.5pt] coordinates {
  (2014,0.6)(2015,0.6)(2016,0.7)(2017,0.9)(2018,1.0)(2019,1.1)
  (2020,2.5)(2021,1.7)(2022,0.9)(2023,0.8)(2024,0.9)
};
\addlegendentry{Vancouver}
\addplot[casered, thick, dashed, mark=square*, mark size=1.5pt] coordinates {
  (2014,1.2)(2015,1.3)(2016,1.1)(2017,1.0)(2018,1.1)(2019,1.1)
  (2020,3.4)(2021,4.3)(2022,1.7)(2023,1.5)(2024,1.5)
};
\addlegendentry{Toronto}
\addplot[dataorange, thick, dotted, mark=triangle*, mark size=1.5pt] coordinates {
  (2014,4.7)(2015,5.3)(2016,7.0)(2017,6.4)(2018,3.8)(2019,3.6)
  (2020,5.2)(2021,4.5)(2022,2.7)(2023,1.4)(2024,1.4)
};
\addlegendentry{Calgary}
\addplot[unbcgreen!80, thick, mark=diamond*, mark size=1.5pt] coordinates {
  (2014,2.7)(2015,2.8)(2016,3.0)(2017,3.0)(2018,2.4)(2019,2.2)
  (2020,3.0)(2021,3.1)(2022,1.9)(2023,1.5)(2024,1.5)
};
\addlegendentry{National average}
\end{axis}
\end{tikzpicture}
\caption{Rental Market Vacancy Rates, Selected Canadian Cities, 2014--2024.
The ``healthy'' vacancy rate of 3\% provides a benchmark for a functioning
rental market. Vancouver has been below 3\% for all but one year since at
least 2014. The COVID-19 pandemic temporarily elevated Toronto's vacancy to
4.3\% in 2021, but this relief was transitory: by 2022, all major markets
had returned to near-zero vacancy. Calgary's trajectory is distinctive ---
the oil price collapse of 2015--2016 created a temporary glut, but the
city subsequently tightened as energy employment recovered and immigration
accelerated post-2021.}
\label{fig:vacancy_rates}
\source{Canada Mortgage and Housing Corporation, Rental Market Report, 2024.}
\end{figure}

When vacancy rates fall below 2\%, the rental market becomes a landlords'
market. Average asking rents for 2-bedroom apartments reached \$3{,}200/month
in Vancouver and \$2{,}750/month in Toronto in late 2024. For a household
earning the provincial median income, these rents consume 40--55\% of
after-tax income.

\begin{figure}[H]
\centering
\begin{tikzpicture}
\begin{axis}[
  ybar,
  bar width=0.40cm,
  width=13cm, height=6.5cm,
  symbolic x coords={{Vancouver},{Toronto},{Ottawa},{Halifax},
                     {Calgary},{Montreal},{Winnipeg},{National Avg}},
  xtick=data,
  x tick label style={rotate=38, anchor=east, font=\small},
  ylabel={Shelter Cost Burden (\% of After-Tax Income)},
  ylabel style={font=\small},
  ymin=0, ymax=72,
  ytick={0,10,20,30,40,50,60,70},
  yticklabel style={font=\small},
  ymajorgrids=true,
  grid style={dashed, gray!25},
  enlarge x limits=0.08,
  legend pos=north east,
  legend style={font=\small},
  clip=false
]
\addplot[fill=casered!60, draw=casered!80] coordinates {
  ({Vancouver},67)({Toronto},58)({Ottawa},46)({Halifax},48)
  ({Calgary},36)({Montreal},42)({Winnipeg},28)({National Avg},50)
};
\addlegendentry{Renter, bottom income quintile}
\addplot[fill=policyblue!50, draw=policyblue!70] coordinates {
  ({Vancouver},48)({Toronto},42)({Ottawa},32)({Halifax},35)
  ({Calgary},24)({Montreal},28)({Winnipeg},22)({National Avg},36)
};
\addlegendentry{Renter, median income}
\addplot[fill=unbcgreen!50, draw=unbcgreen!70] coordinates {
  ({Vancouver},24)({Toronto},21)({Ottawa},17)({Halifax},16)
  ({Calgary},14)({Montreal},15)({Winnipeg},12)({National Avg},18)
};
\addlegendentry{Owner, median income (existing mortgage)}
\addplot[gray!60, dashed, thick, forget plot]
  coordinates {({Vancouver},30)({National Avg},30)};
\node[font=\footnotesize, color=gray!70, align=right]
  at (axis cs:{Montreal},31.5) {30\% threshold};
\end{axis}
\end{tikzpicture}
\caption{Shelter Cost Burden by Tenure and Income Group, Major Canadian Cities, 2024.
Shelter cost burden is annual shelter costs as a percentage of annual after-tax
income. The 30\% threshold (grey dashed line) is the conventional affordability
standard. Low-income renters (bottom quintile) spend over two-thirds of
after-tax income on shelter in Vancouver. Even median-income renters exceed
the 30\% threshold in Vancouver and Toronto. By contrast, median-income
homeowners with existing mortgages --- often locked in at pre-2022 rates ---
face a far smaller shelter cost share, illustrating the growing economic divide
between housing ``haves'' (established owners) and ``have-nots''
(renters and prospective first-time buyers).}
\label{fig:shelter_burden}
\source{Statistics Canada, Survey of Household Spending, 2023--2024;
CMHC Rental Market Report, 2024. Author calculations.}
\end{figure}

%==========================================================
\section{Case Studies}
\label{sec:ch3_casestudies}
%==========================================================

\begin{casestudy}{Vancouver: From Affordable to Global Outlier, 1980--2026}

\textbf{Background:} In 1980, the average house price in Metro Vancouver
was approximately \$120{,}000. Average household income was approximately
\$28{,}000. The price-to-income ratio was 4.3. By 2022, the average price
had reached \$1.37 million while average household income was approximately
\$107{,}000. The PTI had reached 12.8 --- among the highest in the world.

\textbf{What Changed?}

\begin{itemize}
  \item \textbf{Geographic constraints:} Metro Vancouver is squeezed between
        mountains, the US border, the ocean, and the Agricultural Land Reserve.
        The ALR, established in 1973 to protect farmland, constrains the urban
        growth boundary and reduces the supply of developable land.

  \item \textbf{Low supply elasticity:} Single-family zoning covers
        approximately 70\% of Metro Vancouver's residential land. Height limits
        restrict building up; the ALR restricts building out.

  \item \textbf{Capital flows and international demand:} Beginning in the
        late 1980s, Metro Vancouver became a destination for substantial capital
        inflows from Asia --- Hong Kong in the 1990s, mainland China in the
        2000s--2010s. The correlation between capital account openness and
        Vancouver price acceleration is visible in Figure~\ref{fig:house_prices}.

  \item \textbf{Interest rate cycles:} Vancouver prices are highly sensitive
        to real interest rates via the user cost mechanism, amplifying each
        low-rate period.
\end{itemize}

\textbf{Policy Responses:} BC and Vancouver have deployed virtually every
policy tool: 15--20\% Foreign Buyer Tax (2016), BC Speculation and Vacancy
Tax (2018), mortgage stress test (federal, 2018), transit-oriented development
zoning (2023), Small Scale Multi-Unit Housing legislation (2023), and
short-term rental restrictions (2024).

\textbf{Outcome:} Despite these interventions, Vancouver remains among the
three least affordable cities in the world. The 2022--2023 rate-hike correction
reduced prices approximately 14\%, but prices remained 25--30\% above 2019
levels. The fundamental supply constraint has not been resolved. The lesson:
demand-side interventions alone cannot solve a problem rooted in supply
inelasticity and decades of underbuilding.
\end{casestudy}

%----------------------------------------------------------
\begin{casestudy}{Canada vs. Germany: What Different Housing Policies Produce}

\textbf{Background:} Germany provides one of the most instructive international
comparisons. Both countries are wealthy, democratic, urbanised societies with
significant immigrant populations. Yet their housing outcomes differ dramatically.

\textbf{German Housing Outcomes (2024):}
\begin{itemize}
  \item Berlin price-to-income ratio: approximately 8.9
  \item National homeownership rate: approximately 50\% (the lowest in the EU)
  \item Long-term rental tenure: approximately 60\% of households rent,
        many in the same unit for 10--20+ years
  \item Purpose-built rental stock: approximately 55\% of total dwellings
\end{itemize}

\textbf{What Germany Does Differently:}
\begin{enumerate}
  \item \textbf{Tenant protection:} Strong legal protection against eviction
        allows households to plan long-term without fear of displacement,
        reducing the ``premium'' on homeownership as the only path to
        residential stability.

  \item \textbf{Large social housing sector:} Germany's social housing stock
        is 5--6 times larger (as a share of dwellings) than Canada's.

  \item \textbf{Cooperative housing:} Cooperatives own approximately 7\%
        of German housing stock; members pay below-market cost rents and
        have tenure security without speculative price exposure.

  \item \textbf{Transfer taxes on resale:} Germany charges 3.5--6.5\%
        land transfer taxes on property sales, reducing speculative trading.
\end{enumerate}

\textbf{Implication for Canada:} Canada's high homeownership aspiration and
underdeveloped rental sector are not inevitable features of a wealthy economy ---
they are the product of specific policy choices. The German comparison suggests
that sustained, multi-decade commitments to building non-market housing and
protecting tenant tenure can produce meaningfully better affordability outcomes.
\end{casestudy}

%----------------------------------------------------------
\begin{casestudy}{Rent Control: The Economics of a Popular Policy}

\textbf{Background:} Several Canadian provinces impose rent control ---
limits on rent increases for existing tenants. Ontario's rent increase
guideline caps annual increases at the Ontario CPI rate for most units.
BC has similar provisions.

\textbf{The Economic Case Against Rent Control:}
A binding rent ceiling $\bar{R}$ set below market rent $R^*$ creates:
\begin{itemize}
  \item \textbf{Excess demand:} At $\bar{R} < R^*$, quantity demanded
        exceeds quantity supplied. Non-price allocation mechanisms
        (personal connections, waiting lists) replace the price mechanism.
  \item \textbf{Reduced supply:} If developers anticipate future capped
        rents, expected returns on purpose-built rental projects fall,
        reducing new investment. Ontario explicitly exempts units built
        after November 2018 from rent control to address this.
  \item \textbf{Misallocation:} Tenants have strong incentives to remain
        in controlled units even when their circumstances change, reducing
        labour market mobility.
  \item \textbf{Reduced maintenance:} Landlords facing controlled rents
        may reduce maintenance spending to cut costs.
\end{itemize}

\textbf{The Economic Case For Rent Control:}
When vacancy rates are near zero, landlords have market power and can
extract rents above competitive levels. Rent control can correct a market
power problem, provide tenure security, and prevent displacement of
established low-income communities.

\textbf{The Evidence:}
A 2019 Stanford study found that San Francisco's rent control reduced
rental housing supply by 15\% as landlords converted units to condos
or redeveloped to escape controls. Ontario's pattern of low purpose-built
rental construction throughout the 1990s--2000s (when rent control was
broadly applied), followed by a surge in purpose-built starts after the
2018 exemption, is consistent with the supply-reduction effect.

\textbf{The Policy Synthesis:} Rent control protects \textit{existing}
tenants at the expense of \textit{future} tenants and new supply. Where
tenant protection and tenure security are the goals, alternatives such
as vacancy controls, tenant right-of-first-refusal, and anti-eviction
legislation may achieve those goals with less supply distortion. Where
affordability is the goal, housing allowances and social housing expand
supply rather than constrain it.
\end{casestudy}

%==========================================================
\section{Social Housing: Canada in International Context}
\label{sec:social_housing}
%==========================================================

\begin{figure}[H]
\centering
\begin{tikzpicture}
\begin{axis}[
  ybar,
  bar width=0.50cm,
  width=13cm, height=6.5cm,
  symbolic x coords={{USA},{Canada},{Australia},{Germany},
                     {France},{UK},{Denmark},{Netherlands}},
  xtick=data,
  x tick label style={font=\small},
  ylabel={Social Housing Stock (\% of Total Dwellings)},
  ylabel style={font=\small},
  ymin=0, ymax=35,
  ytick={0,5,10,15,20,25,30},
  yticklabel style={font=\small},
  ymajorgrids=true,
  grid style={dashed, gray!25},
  enlarge x limits=0.10,
  nodes near coords,
  nodes near coords style={font=\small},
  every axis plot/.append style={fill=exampleteal!60, draw=exampleteal!80}
]
\addplot coordinates {
  ({USA},     1.0)
  ({Canada},  3.5)
  ({Australia},4.8)
  ({Germany}, 12.5)
  ({France},  17.0)
  ({UK},      17.5)
  ({Denmark}, 21.0)
  ({Netherlands},30.5)
};
\end{axis}
\end{tikzpicture}
\caption{Social Housing as a Share of Total Dwellings, Selected Countries.
Canada's social housing stock (approximately 3.5\%) is among the smallest
of comparable wealthy nations. The Netherlands (30.5\%) and Denmark (21\%)
maintain large non-market housing sectors that insulate a significant share
of their populations from private market volatility. The OECD recommends
that countries targeting housing affordability increase non-market supply
to at least 10--15\% of total dwellings. Achieving this in Canada would
require building approximately 750{,}000--1{,}000{,}000 additional social
and affordable units --- a multi-decade project requiring sustained political
commitment and federal investment.}
\label{fig:social_housing}
\source{OECD Affordable Housing Database, 2023; CMHC National Housing Strategy, 2024.}
\end{figure}

%==========================================================
\section{Federal Housing Policy: What Has Been Tried}
\label{sec:housing_policy}
%==========================================================

\begin{table}[H]
\centering
\caption{Major Federal Housing Policy Measures, 2015--2026}
\label{tab:housing_policies}
\renewcommand{\arraystretch}{1.30}
\begin{tabular}{p{3.3cm}p{1.0cm}p{2.0cm}p{4.6cm}p{2.6cm}}
\toprule
\textbf{Policy} & \textbf{Year} & \textbf{Type} & \textbf{Mechanism}
                & \textbf{Assessment} \\
\midrule
Mortgage stress test
  & 2018 & Demand (regulatory)
  & Borrowers qualify at contract rate + 2\%, reducing borrowing capacity
  & Effective at limiting leverage; worsens access for marginal buyers \\
\addlinespace
First Home Savings Account
  & 2023 & Demand (tax)
  & \$40K tax-free savings for first-time buyers; contributions tax-deductible
  & Modest benefit; may capitalise into higher prices \\
\addlinespace
FTHBI (Shared Equity)
  & 2019 & Demand (subsidy)
  & CMHC takes 5--10\% equity stake; disbanded 2024
  & Limited uptake; demand subsidy in supply-constrained market \\
\addlinespace
National Housing Co-Investment Fund
  & 2018 & Supply (subsidy)
  & Low-interest loans for affordable rental construction
  & Positive but underfunded; units produced below target \\
\addlinespace
GST Removal on Purpose-Built Rentals
  & 2023 & Supply (tax)
  & Eliminates 5\% GST on new purpose-built rental construction
  & Positive; reduces effective construction cost by 3--5\% \\
\addlinespace
Housing Accelerator Fund
  & 2023 & Supply (incentive)
  & \$4B fund to municipalities contingent on zoning reform
  & Promising; targets root constraint; early results positive \\
\addlinespace
Foreign Buyer Ban
  & 2023 & Demand (regulatory)
  & Prohibits non-residents from buying residential property
  & Marginal effect; foreign buyers were small share of transactions \\
\bottomrule
\end{tabular}
\source{Department of Finance Canada; CMHC; Parliamentary Budget Officer,
Housing Affordability Assessment, 2024.}
\end{table}

The dominant pattern in Table~\ref{tab:housing_policies} is instructive:
most federal housing policies through 2022 were \textbf{demand-side}
measures --- tax credits, subsidies, and regulatory restrictions. With a
near-vertical supply curve (Panel A of Figure~\ref{fig:housing_sd}),
demand-side interventions cannot improve affordability. A demand subsidy
to a fixed-supply market raises prices; a demand restriction may lower
prices briefly but also reduces ownership opportunity for the households
the policy was meant to help.

The shift from 2022 onwards toward supply-side measures --- the Housing
Accelerator Fund, the GST removal on purpose-built rentals, federal pressure
on municipalities to upzone --- reflects a belated recognition of this basic
supply-and-demand reality.

%==========================================================
\begin{policydebate}{Is Foreign Ownership Driving Canadian Housing Unaffordability?}

\textbf{The Issue:} Few housing policy questions generate more political
heat and less empirical consensus than the role of foreign ownership. Advocates
argue that offshore capital has driven Vancouver and Toronto prices beyond
what domestic incomes could support. Sceptics argue foreign ownership is a
convenient scapegoat for domestic policy failures.

\textbf{Evidence Supporting a Significant Foreign Effect:}
\begin{itemize}
  \item CMHC's non-resident ownership data showed that non-residents owned
        approximately 2.4\% of residential properties in Ontario and 4.6\%
        in British Columbia, with concentrations in condominium markets.
  \item The 2016 BC Foreign Buyer Tax caused an immediate and dramatic
        (though temporary) cooling of Vancouver prices, particularly in
        the detached segment most favoured by offshore buyers.
\end{itemize}

\textbf{Evidence Against a Large Foreign Effect:}
\begin{itemize}
  \item Aggregate non-resident ownership of 2--5\% of total housing
        stock is insufficient to explain prices 3--5 times higher
        than domestic income fundamentals alone would support. If foreign
        ownership were the primary driver, the 2023 foreign buyer ban
        should have made housing affordable --- it did not.
  \item Canadian housing has appreciated across \textit{all} markets,
        including cities (Halifax, Saskatoon, Prince George) with essentially
        zero foreign ownership activity. A supply-constrained explanation
        applies everywhere; a foreign-buyer explanation applies only to
        a few markets.
  \item Many studies find that domestic investors and domestic demographic
        pressure are quantitatively more important drivers.
\end{itemize}

\textbf{Policy Implication:} The evidence is clear: Canada cannot tax its
way to housing affordability. Foreign buyer restrictions are not irrelevant,
but they address a marginal contributor to a problem driven primarily by
supply constraints and domestic demand. Supply must grow.
\end{policydebate}

%==========================================================
\begin{everydaylife}{The Mathematics of Your First Home}

Suppose you graduate from UNBC at age 22 earning \$58{,}000 per year in
Prince George, with student debt of \$28{,}000. The current benchmark price
for a two-bedroom home in Prince George is approximately \$380{,}000.

\textbf{Down payment:} A 5\% minimum down payment requires \$19{,}000,
plus CMHC insurance of approximately \$7{,}600 added to your loan.
A 20\% down payment (\$76{,}000) eliminates the insurance.

\textbf{Saving timeline:} After tax, rent (\$1{,}500/month), loan payments
(\$400/month), and living costs (\$1{,}800/month), you might save
\$600--\$1{,}000/month. At \$800/month, you accumulate \$76{,}000 in about
8 years --- but if prices appreciate at 3\% annually while you save, the
20\% down payment target grows from \$76{,}000 to \$96{,}200. You are on
a treadmill.

\textbf{At purchase (age 30):} A 25-year mortgage on a \$304{,}000 principal
at 4.5\% costs approximately \$1{,}655/month --- 28.6\% of your gross monthly
income of \$5{,}800. This is within the 30\% guideline, but tight.

\textbf{Key insight from the user cost formula:} Whether buying is rational
depends on the gap between $r$ (your mortgage rate) and $g$ (expected future
appreciation). If prices are flat or falling, the user cost rises sharply.
If rates fall and prices appreciate, the user cost makes owning attractive.
You cannot control either variable --- they are set by the Bank of Canada
and the housing market. What you can control is your savings rate and your
First Home Savings Account contributions. The mathematics of homeownership
in modern Canada is more uncertain, more leveraged, and more dependent on
external forces than at any previous point in the country's history.
\end{everydaylife}

%==========================================================
\section*{Chapter Summary}
%==========================================================
\addcontentsline{toc}{section}{Chapter Summary}

\begin{itemize}
  \item Housing is simultaneously a consumption good, a durable asset,
        a financial investment, and a location-specific good. These
        characteristics make housing markets unusually sensitive to
        interest rates, supply constraints, and demographic shifts.

  \item The short-run housing supply curve is nearly vertical in Canadian
        cities (supply elasticity $\varepsilon_s \approx 0.2$--$0.6$) due
        to zoning restrictions, slow permitting, development charges, and
        construction labour shortages. When supply is inelastic, demand
        shocks produce large price increases with minimal quantity response.

  \item The \textbf{user cost formula} ($\text{UC} = (r + \delta + \tau - g) \times P$)
        explains the rent-vs-buy decision. In 2021, deeply negative real
        rates and high expected capital gains drove user costs near zero,
        generating enormous ownership demand. By 2023, user costs had risen
        six-fold, making renting more economical in many markets.

  \item The \textbf{mortgage payment formula} (Equation~\ref{eq:mortgage_payment})
        shows that at 5.8\% for a \$700{,}000 home with 20\% down,
        the monthly payment is \$3{,}531, requiring gross income of at
        least \$141{,}240 to meet the 30\% affordability threshold.

  \item Vancouver (PTI 12.8) and Toronto (PTI 10.2) rank among the least
        affordable cities globally. The CMHC affordability index for
        Vancouver reached 92.1\% in 2024.

  \item House prices have risen more than four-fold nationally since 2000.
        The 2020--2022 COVID-era rate cuts triggered the most explosive
        price surge in Canadian history. The 2022--2023 rate hike cycle
        produced the first significant correction but left prices far
        above pre-pandemic levels.

  \item The supply-demand gap is structural: Canada builds 240{,}000--265{,}000
        units per year but needs 400{,}000--480{,}000 at current immigration
        levels. This shortfall drives both ownership price appreciation
        and the rental vacancy crisis.

  \item The rental vacancy rate has fallen below 2\% in most major
        Canadian cities, producing double-digit rent inflation and
        acute affordability problems for low-income renters.

  \item Most federal housing policies through 2022 were demand-side.
        Economic analysis demonstrates that demand subsidies in
        supply-constrained markets raise prices rather than improving
        affordability. Post-2022 policy has shifted toward supply-side
        instruments.

  \item International comparisons show that large non-market housing
        sectors, strong tenant protections, and permissive zoning produce
        significantly better affordability outcomes in comparable democracies.
\end{itemize}

%==========================================================
\section*{Key Terms}
%==========================================================
\addcontentsline{toc}{section}{Key Terms}

\begin{description}[style=nextline, leftmargin=3.5cm]
  \item[\term{Development charges}] Fees levied by municipalities on
        new construction to fund infrastructure costs.
  \item[\term{Financialization}] The increasing treatment of residential
        real estate as a financial asset rather than primarily as shelter.
  \item[\term{Housing starts}] New residential units on which construction
        has begun in a given period.
  \item[\term{Housing supply elasticity}] $\varepsilon_s = \%\Delta Q_s / \%\Delta P$;
        how responsive new construction is to house price changes.
  \item[\term{Normal good}] A good for which demand rises with income;
        housing has income elasticity of approximately 0.5--1.0.
  \item[\term{Mortgage stress test}] Federal requirement to qualify at
        contract rate + 2\% or the Bank of Canada benchmark rate,
        whichever is higher.
  \item[\term{Price-to-income ratio (PTI)}] Median house price divided by
        median annual household income; the primary affordability benchmark.
  \item[\term{Purpose-built rental}] Buildings constructed specifically
        for the long-term rental market.
  \item[\term{Rent control}] Government regulations limiting annual
        rent increases.
  \item[\term{Rental vacancy rate}] Percentage of purpose-built rental
        units that are unoccupied and available.
  \item[\term{Shelter cost burden}] Shelter costs as a share of household
        income; above 30\% is conventionally ``unaffordable.''
  \item[\term{Social housing}] Housing owned by government or non-profits
        and rented at below-market rates to income-qualifying households.
  \item[\term{Upzoning}] A land-use regulation change permitting higher-
        density development than was previously allowed.
  \item[\term{User cost of housing}] Annual economic cost of owning a home:
        $\text{UC} = (r + \delta + \tau - g) \times P$.
  \item[\term{Vacancy tax}] A tax on vacant residential properties,
        incentivising owners to rent units out.
  \item[\term{Zoning}] Municipal regulation specifying permitted land uses
        and development standards.
\end{description}

%==========================================================
\section*{Review Questions}
%==========================================================
\addcontentsline{toc}{section}{Review Questions}

\begin{enumerate}

\item \textbf{(Concept)} Explain why the short-run housing supply curve
is nearly vertical in major Canadian cities. What are the three most
important constraints on rapid housing supply expansion? Which is most
amenable to policy reform?

\item \textbf{(Application)} Using the user cost formula
(Equation~\ref{eq:user_cost}), calculate the annual user cost of owning
a \$900{,}000 home in each scenario. In which is buying most attractive
relative to paying \$2{,}800/month rent?
\begin{enumerate}[(a)]
  \item $r = 2\%$, $\delta = 2\%$, $\tau = 0.8\%$, $g = 5\%$
  \item $r = 5.5\%$, $\delta = 2\%$, $\tau = 0.8\%$, $g = 1\%$
  \item $r = 4.0\%$, $\delta = 2\%$, $\tau = 0.8\%$, $g = 3\%$
\end{enumerate}

\item \textbf{(Measurement)} Canada's average house price rose from
approximately \$530{,}000 in 2020 to \$796{,}000 in 2022.
\begin{enumerate}[(a)]
  \item What was the nominal percentage price increase?
  \item If median household income rose from \$84{,}000 to \$90{,}000
        over the same period, what happened to the PTI?
  \item What does the PTI change tell us that the nominal price increase alone does not?
\end{enumerate}

\item \textbf{(Policy)} A policymaker proposes a \$30{,}000 grant to
all first-time homebuyers. Using Figure~\ref{fig:housing_sd}, explain
what happens to house prices in: (a) a market with highly inelastic
supply; (b) a market with elastic supply. Which scenario reflects most
Canadian urban markets, and what does this imply about the policy's
effectiveness?

\item \textbf{(Rent control)} Using supply-and-demand analysis, explain
the short-run and long-run effects of a binding rent ceiling set below
the market rent. What is the key difference between the short-run and
long-run outcomes, and why does this matter for policy evaluation?

\item \textbf{(International comparison)} Canada's social housing stock
is approximately 3.5\% of total dwellings; the Netherlands is 30.5\%.
Explain the mechanism by which a larger social housing sector would
affect equilibrium rents in the private market. What are two obstacles
to Canada rapidly expanding its social housing stock?

\end{enumerate}

%==========================================================
\section*{Quantitative Problems}
%==========================================================
\addcontentsline{toc}{section}{Quantitative Problems}

\begin{enumerate}

\item \textbf{Mortgage Payment Calculations.}
Use Equation~\ref{eq:mortgage_payment} with Canadian semi-annual compounding
($r_m = (1 + r_{\text{annual}}/2)^{1/6} - 1$) and 25-year amortization.
\begin{enumerate}[(a)]
  \item A \$600{,}000 home, 10\% down, at 4.0\% annual rate. Monthly payment?
  \item Same home with 20\% down at 4.0\%. Monthly payment?
  \item Same home with 20\% down at 6.5\%. Monthly payment?
  \item For each scenario, what gross annual income is required for the
        payment to equal exactly 30\% of monthly gross income?
\end{enumerate}

\item \textbf{User Cost Analysis.}
A household in Kelowna, BC is choosing between renting and buying.
The home costs \$650{,}000; a comparable rental unit rents for
\$2{,}400/month (\$28{,}800/year).
\begin{enumerate}[(a)]
  \item Calculate the user cost under 2021 conditions ($r=2\%$, $\delta=2\%$,
        $\tau=0.8\%$, $g=6\%$) and under 2023 conditions ($r=5.8\%$,
        $\delta=2\%$, $\tau=0.8\%$, $g=0\%$). Is renting or buying more
        economical in each year?
  \item What annual capital gain rate $g$ (holding 2023 parameters constant)
        would make the household indifferent between renting and buying?
        Show your algebra.
  \item If the home price falls to \$560{,}000 but the rate remains at
        5.8\%, calculate the new user cost. Has affordability improved?
\end{enumerate}

\item \textbf{Housing Supply Elasticity.}
In Market A (Toronto), a 20\% price increase led to a 4\% increase in
housing starts. In Market B (Houston), the same 20\% increase led to
a 30\% increase in starts.
\begin{enumerate}[(a)]
  \item Calculate the price elasticity of housing supply in each market.
  \item Draw the supply curves, starting from common initial conditions
        of \$400{,}000 and 50{,}000 units.
  \item Suppose demand increases by 30\% in both markets (demand
        elasticity = $-0.8$). What is the new equilibrium price in each?
  \item What policy changes would move Market A toward Market B's
        supply elasticity?
\end{enumerate}

\item \textbf{Affordability Index Calculation.}
A city has the following characteristics: median benchmark price
\$820{,}000; 20\% down payment (\$164{,}000); 5.0\% mortgage rate;
25-year amortization; 0.65\% annual property tax rate; \$250/month
utilities; median annual household income \$102{,}000.
\begin{enumerate}[(a)]
  \item Calculate the monthly mortgage payment.
  \item Calculate the monthly property tax payment.
  \item Calculate the CMHC Affordability Index (total monthly housing
        costs as \% of median monthly gross income).
  \item Compare to Table~\ref{tab:affordability_index}. Which Canadian
        city does this most closely resemble?
\end{enumerate}

\item \textbf{Supply-Demand Gap.}
Canada's population grew by 1{,}053{,}000 persons in 2022.
\begin{enumerate}[(a)]
  \item At an average household size of 2.5 persons, how many new
        households were formed in 2022?
  \item Adding a demolition rate of 0.5\% of the existing stock
        ($\approx 75{,}000$ units), what total units were needed?
  \item Actual starts were 263{,}000. What was the shortfall?
  \item If this shortfall persists for 5 years, how large a housing
        deficit accumulates? What effect would you expect on rents?
\end{enumerate}

\item \textbf{Data Exercise.} Visit the CMHC website
(\url{www.cmhc-schl.gc.ca}) and find the most recent Housing Market
Outlook for your city.
\begin{enumerate}[(a)]
  \item Report the current rental vacancy rate and year-over-year
        change in average asking rents.
  \item Report housing starts over the past 12 months.
  \item Compare starts to CMHC's estimated required units. Is supply
        growing faster or slower than demand?
  \item Identify one supply-side policy your municipal or provincial
        government has announced. Using Figure~\ref{fig:housing_sd},
        evaluate whether it targets the right constraint.
\end{enumerate}

\end{enumerate}

%==========================================================
\section*{Essay Questions}
%==========================================================
\addcontentsline{toc}{section}{Essay Questions}

\begin{enumerate}

\item \textbf{(600--800 words)} ``Canada's housing crisis is fundamentally
a supply problem, and demand-side policies cannot solve it.'' Evaluate this
claim using the supply-and-demand model. Your essay should: (a) explain
why the short-run supply curve is nearly vertical; (b) use a diagram to
show why demand subsidies raise prices in supply-constrained markets;
(c) describe at least two supply-side policies and evaluate their
effectiveness; and (d) discuss whether demand-side interventions could
complement effective supply-side reforms.

\item \textbf{(500--700 words)} Compare Canada's housing outcomes to those
of Germany (Case Study 2). What are the three most important institutional
differences between the two countries? For each, assess the feasibility and
political obstacles to implementing a similar approach in Canada. Would you
recommend that Canada move toward the German model? Defend your answer.

\item \textbf{(400--600 words)} The user cost of housing in 2021 was near
zero in many Canadian markets, while the same calculation in 2023 yielded
a cost above 8\% of home value. Write an essay explaining what drove this
change, what its economic consequences were for the housing market, and
what it tells us about the relationship between monetary policy and housing
affordability. Connect your analysis explicitly to Chapter~\ref{ch:inflation}'s
discussion of the Bank of Canada's rate decisions.

\end{enumerate}

%==========================================================
\section*{Discussion Questions}
%==========================================================
\addcontentsline{toc}{section}{Discussion Questions}

\begin{enumerate}

\item The opening vignette describes three households: Vancouver engineers
who bought for \$2.1M, a nurse and teacher who cannot afford to buy, and a
Prince George mill worker paying high rent. Using concepts from this chapter,
explain why these three households face such different situations despite
the same national macroeconomic policies. Which economic force is most
responsible for each household's situation?

\item BC's 2023 upzoning legislation allows multiplexes on formerly
single-family lots. Who wins and who loses from this policy? Consider
existing homeowners, would-be first-time buyers, renters, developers,
and local municipalities. Does the aggregate economic benefit outweigh
the distributional costs?

\item Should Canada maintain the Agricultural Land Reserve in BC, given
that it constrains the supply of developable land around Metro Vancouver
and contributes to housing unaffordability? Or does food security and
environmental preservation justify the housing cost? How would an economist
structure the argument for each side?

\item The Case Study on rent control concludes that it protects existing
tenants at the expense of future tenants and new supply. Is this an
acceptable trade-off? Consider the perspective of: (a) a long-term renter
in a controlled unit in Toronto; (b) a newcomer searching for an apartment
in the same city; and (c) a developer deciding whether to build new rental
units. Is there a policy that achieves rent control's goals without its
supply-side costs?

\end{enumerate}
